\ifdefined\iscombined
	\tutorial{Tutorial 1}
\else
\documentclass{article}
\usepackage{xparse}
\usepackage{tabularx}
\usepackage{changepage}
\usepackage{listings}
\usepackage{amsthm}
\usepackage{comment}
\usepackage{amsmath}
\usepackage{braket}
\usepackage{amsfonts}
\usepackage{amssymb}
\usepackage{sectsty}
\usepackage{hyperref}
\usepackage{fancyhdr}
\usepackage[top=1in,bottom=1in,left=1in,right=1in]{geometry}
\usepackage{float}
\usepackage{graphicx}
\usepackage{mathtools}
\usepackage{tikz}
\usetikzlibrary{arrows}

\date{
	\vspace{-.75em}
	{\large \today}
}

\title{
	\vspace*{-1.5em}
	{\Huge Tutorial 1} \\
	\vspace{-1em}
	\rule{\textwidth/2}{.01em} \\
	\vspace{-.75em}
}

\author{
	{\large Matthew Farah}
}

\hypersetup{
	colorlinks=true,
	linkcolor=blue,
	filecolor=magenta,
	urlcolor=blue,
	citecolor=blue,
	anchorcolor=blue
}

\fancyhead{}
\fancyhead[L]{\today}
\fancyhead[R]{Tutorial 1}

\setlength{\parindent}{0cm}

\tikzset{vertex/.style = {shape = circle, draw, minimum size = 2em}}
\tikzset{edge/.style = {->,> = latex'}}

\newenvironment{solution}{
	\vspace{.5em}
	\par
	\begin{adjustwidth}{1.5em}{}
	\textbf{{Solution:}}
	\par
	\vspace{.5em}
}{
	\vspace{1em}
	\end{adjustwidth}
	\setcounter{step}{0}
}

\makeatother
\makeatletter

\newcounter{step}
\renewcommand{\thestep}{\Roman{step}}

\newenvironment{step}{
	\refstepcounter{step}
	\vspace{1em}\par\noindent
	\ignorespaces\noindent\textbf{(\thestep)}
	\ignorespaces
}{
	\par
	\vspace{1.5em}
}

\newcounter{problem}
\renewcommand{\theproblem}{\arabic{problem}}

\newenvironment{problem}[1][]{
	\newpage
	\refstepcounter{problem}
	\setcounter{subproblem}{0}
	\vspace{.5em}\par\noindent
	\begin{adjustwidth}{1.5em}{}
	\noindent\textbf{Problem \theproblem: }\noindent\textit{#1}
	\par
	\phantomsection
	\addcontentsline{toc}{section}{\protect{\large{Problem \theproblem}}}
}{
	\vspace{.5em}
	\end{adjustwidth}
}


\newcounter{subproblem}
\renewcommand{\thesubproblem}{\alph{subproblem}}

\newenvironment{subproblem}{
	\setcounter{step}{0}
	\refstepcounter{subproblem}
	\phantomsection
	\addcontentsline{toc}{subsection}{\protect{\large{(\thesubproblem)}}}
	\begin{adjustwidth}{1.5em}{}
	\noindent{\textbf{(\thesubproblem) }}\ignorespaces
}{
	\par
	\vspace{.5em}
	\end{adjustwidth}
}

\newcounter{lemma}
\renewcommand{\thelemma}{\arabic{lemma}}

\newenvironment{lemma}[1][]{
	\refstepcounter{lemma}
	\vspace{1em}\par\noindent
	\noindent\textbf{Lemma ({#1}): }\noindent\ignorespaces
}{
	\vspace{.5em}
}

\newcounter{proposition}
\renewcommand{\theproposition}{\arabic{proposition}}

\newenvironment{proposition}{
	\refstepcounter{proposition}
	\vspace{1em}\par\noindent
	\noindent\textbf{Proposition:}
	\noindent\textit
	\ignorespaces
}{
	\vspace{.5em}
}

\newcommand{\supp}[1]{\text{supp}({#1})}
\newcommand{\ord}[1]{\text{ord}({#1})}
\newcommand{\gen}[1]{\langle {#1} \rangle}
\newcommand{\lcm}[2]{\text{lcm}({#1}, {#2})}
\renewcommand{\mod}[1]{\ (\text{mod} \; {#1})}

\sectionfont{\fontsize{12}{12}\bfseries}
\subsectionfont{\fontsize{10}{12}\normalfont}
\subsubsectionfont{\fontsize{10}{12}\normalfont}

\begin{document}

\maketitle
\pagestyle{fancy}

\newcolumntype{Y}{>{\centering\arraybackslash}X}

\tableofcontents

\fi

\begin{problem}
	Define a binary operator $\ast$ on $\mathbb{Z} \times \mathbb{Z}$ by:

	\begin{align*}
		(a, b) \ast (c, d) \coloneq (ab + bc, bd) \\
	\end{align*}

	Is $\mathbb{Z} \times \mathbb{Z}$ a group under this operation? If so, prove it by establishing which group axioms fail (with proofs/counterexamples).
\end{problem}

\begin{solution}
	$\mathbb{Z} \times \mathbb{Z}$ does \underline{not} constitute a group under $\ast$. To see why, we'll examine which of the group axioms fail to hold and why. As a refresher, for $(\mathbb{Z} \times \mathbb{Z}, \ast)$ to be a group, it must satisfy all of the following four criteria:
	\begin{enumerate}
		\item $\ast$ must be a binary operator on $\mathbb{Z} \times \mathbb{Z}$.
		\item $\ast$ must be associative over $\mathbb{Z} \times \mathbb{Z}$.
		\item $\ast$'s identity element must be contained in $\mathbb{Z} \times \mathbb{Z}$.
		\item Every element in $\mathbb{Z} \times \mathbb{Z}$ must have an inverse in $\mathbb{Z} \times \mathbb{Z}$ with respect to $\ast$. \\
	\end{enumerate}

	We will investigate each criterion separately below.

	\begin{step}
		Criterion one.
	\end{step}

	The first criterion states that $\ast$ must be a binary operator on $\mathbb{Z} \times \mathbb{Z}$. \\

	For $\ast$ to form a binary operator over $\mathbb{Z} \times \mathbb{Z}$, it must be defined and closer over all of $\mathbb{Z} \times \mathbb{Z}$. By inspection, and knowledge of the fact that $\mathbb{Z}$ is closed under both multiplication and addition, we can see that this is actually the case, and so $\ast$ does indeed constitute a valid binary operator over $\mathbb{Z} \times \mathbb{Z}$.

	\begin{step}
		Criterion two.
	\end{step}

	Criterion two states that $\ast$ must be associative over $\mathbb{Z} \times \mathbb{Z}$. \\

	By counter example, we can see that this fails to hold true. Although $(1, 0), (0, 2)$, and $(3, 0)$ all belong to $\mathbb{Z} \times \mathbb{Z}$, we have that:

	\begin{align*}
		((1, 0) \ast (0, 2)) \ast (3, 0) &= (2, 0) \ast (3, 0) \\
		&= (0, 0) \\
	\end{align*}

	Whereas:

	\begin{align*}
		(1, 0) \ast ((0, 2) \ast (3, 0)) &= (1, 0) \ast (0, 6) \\
		&= (6, 0) \\
	\end{align*}

	And thus, since $(a \ast b) \ast c \neq a \ast (b \ast c)$ for all $a, b, c \in (\mathbb{Z} \times \mathbb{Z})$, $\ast$ is not associative over its domain and so criterion two isn't satisfied.

	\begin{step}
		Criterion three.
	\end{step}

	Criterion three states that $\ast$ must have an identity element in $\mathbb{Z} \times \mathbb{Z}$. \\

	Despite not being associative, it is actually the case that $\ast$ has an identity element in $\mathbb{Z} \times \mathbb{Z}$. To see this, let $(a, b)$ be an arbitrary element in $\mathbb{Z} \times \mathbb{Z}$. Since:

	\begin{align*}
		\boxed{(a, b) \ast (0, 1)} &= (a(1) + b(0), b(1)) &&\text{(By definition of $\ast$)} \\
		&= \boxed{(a, b)} &&\text{(Simplifying)} \\
		&= (0(b) + 1(a), 1(b)) &&\text{(Expanding)} \\
		&= \boxed{(0, 1) \ast (a, b)} &&\text{(By definition of $\ast$)} \\
	\end{align*}

	$(0, 1) \in (\mathbb{Z} \times \mathbb{Z})$ therefore satisfies $(a, b) \ast (0, 1) = (a, b) = (0, 1) \ast (a, b)$ for arbitrary $(a, b) \in \mathbb{Z} \times \mathbb{Z}$, making it an identity element for $\ast$ in its domain. Thus, $\ast$ really does have an identity element in $\mathbb{Z} \times \mathbb{Z}$, satisfying criterion three.

	\begin{step}
		Criterion four.
	\end{step}

	Lastly, criterion four requires that every element in $\mathbb{Z} \times \mathbb{Z}$ have an inverse under $\ast$ that's contained in $\mathbb{Z} \times \mathbb{Z}$. \\

	$\ast$ fails to satisfy criterion four as well. Taking $(1, 0)$ as an example, we can see it to have an inverse in $\mathbb{Z} \times \mathbb{Z}$ there must exist some $(a, b) \in (\mathbb{Z} \times \mathbb{Z})$ such that:

	\begin{align*}
		(1, 0) \ast (a, b) &= (0, 1) &&\text{(By definition of inverse)} \\
	\end{align*}

	Since we determined as a part of our analysis of criterion three that $(0, 1)$ must be the identity of $\ast$ over $\mathbb{Z} \times \mathbb{Z}$. However, since:

	\begin{align*}
		(1, 0) \ast (a, b) &= (0, 1) &&\text{(By our earlier equality)} \\
		\implies ((1)(b) + (0)(a), (0)(b)) &= (0, 1) &&\text{(By definition of $\ast$)} \\
		\implies (b, 0) &= (0, 1) &&\text{(Simplifying)} \\
		\implies 0 &= 1 &&\text{(Equating entries)} \\
	\end{align*}

	Since this clearly cannot be the case, $(1, 0)$ (and actually, any element of the form $(x, 0)$) cannot have an inverse in $\mathbb{Z} \times \mathbb{Z}$ under $\ast$. Hence, criterion four is violated as well.
\end{solution}

\begin{problem}
	Let $G$ be a group and let $a, b \in G$. Prove that the inverse of $ab$ is $b^{-1}a^{-1}$.
\end{problem}

\begin{solution}
	Let $a, b \in G$ be arbitrary and let $a^{-1}, b^{-1} \in G$ be their respective inverses. To show that the inverse of $ab$ is $b^{-1}a^{-1}$, we need to show that $(ab)(b^{-1}a^{-1}) = (b^{-1}a^{-1})(ab) = 1$. Thankfully, this can be proven directly and relatively easily using the following manipulations:

	\begin{align*}
		\boxed{(ab)(b^{-1}a^{-1})} &= a(b(b^{-1}a^{-1})) &&\text{(By associativity)} \\
		&= a((bb^{-1})a^{-1}) &&\text{(By associativity)} \\
		&= a((1)a^{-1}) &&\text{(By definition of inverse)} \\
		&= aa^{-1} &&\text{(By definition of identity)} \\
		&= \boxed{1} &&\text{(By definition of inverse)} \\
		&= b^{-1}b &&\text{(By definition of inverse)} \\
		&= b^{-1}((1)b) &&\text{(By definition of identity)} \\
		&= b^{-1}((a^{-1}a)b) &&\text{(By definition of inverse)} \\
		&= (b^{-1}(a^{-1}a))b &&\text{(By associativity)} \\
		&= ((b^{-1}a^{-1})a)b &&\text{(By associativity)} \\
		&= \boxed{(b^{-1}a^{-1})(ab)} &&\text{(By associativity)} \\
	\end{align*}

	And so, the inverse of $ab$ is indeed $b^{-1}a^{-1}$!
\end{solution}

\begin{problem}
	Let $G$ be a group and let $a, b \in G$. Prove that if $a$ and $b$ commute, then so do $a^{-1}$ and $b^{-1}$.
\end{problem}

\begin{solution}
	Once again, let $a, b \in G$ be arbitrary elements such that $a$ and $b$ commute, and let $a^{-1}, b^{-1} \in G$ be their respective inverses. To show that $a^{-1}$ and $b^{-1}$ commute, we must prove that $a^{-1}b^{-1} = b^{-1}a^{-1}$. Since $a$ and $b$ are assumed to be commutative, we know that $ab = ba$. Using this, we have that: \\

	\begin{align*}
		a^{-1}b^{-1} &= 1(a^{-1}b^{-1}) &&\text{(By definition of identity)} \\
		&= ((b^{-1}a^{-1})(ab))(a^{-1}b^{-1}) &&\text{(By definition of inverse)} \\
		&= (b^{-1}a^{-1})((ab)(a^{-1}b^{-1})) &&\text{(By associativity)} \\
		&= (b^{-1}a^{-1})((ba)(a^{-1}b^{-1})) &&\text{(By commutativity of $a$ and $b$)} \\
		&= (b^{-1}a^{-1})(1) &&\text{(By definition of inverse)} \\
		&= b^{-1}a^{-1} &&\text{(By definition of identity)} \\
	\end{align*}

	And so, since $a^{-1}b^{-1} = b^{-1}a^{-1}$, we indeed have that $a^{-1}$ and $b^{-1}$ commute.
\end{solution}

\begin{problem}
	Prove that a group with exactly four elements must be abelian. \\

	(\emph{Hint}: Let $G \coloneq \set{1, a, b, c}$ and make a multiplication table for $G$ by brute force. Start by showing either that $ab = 1$ or $ab = c$; thus proceed by brute force.)
\end{problem}

\begin{solution}

\end{solution}

\ifdefined\iscombined
	\newpage
\else
	\end{document}
\fi
